%% Chapter 4 -----------------------------------------------------------------
\chapter{Help, Updates and Support}
\label{ch:help}

This chapter describes the built-in user guide, the citation dialog, the update check, the
bug-report tool and the version information.

\section{The built-in user guide}
\label{sec:in-app-guide}
\index{user guide!built-in}

\menu{Help > User Guide} (\kbd{F1}) opens a condensed version of this guide in a separate
window (Figure~\ref{fig:in-app-guide}). The guide follows the application theme and offers:
\begin{itemize}
  \item a linked table of contents and cross-references between sections;
  \item \gui{Home} and \gui{Back} buttons for navigation;
  \item a \gui{Find in guide...} field with \gui{Find Next}; and
  \item \gui{Close}.
\end{itemize}
The GCS CME Fitting window has its own help, opened with \kbd{F1} in that window
(Section~\ref{sec:gcs-menus}).

\screenshot[0.7\linewidth]{dialogs/user_guide_dialog}{The built-in user guide window.}{fig:in-app-guide}{The \emph{e-CALLISTO FITS Analyzer User Guide} window opened with F1, showing the table of contents, the Home/Back buttons and the Find field.}

\section{Citing the software}
\label{sec:cite-dialog}
\index{citation!dialog}

The \gui{Cite this Software}\index{Cite this Software} button at the right end of the toolbar opens the citation
dialog (Figure~\ref{fig:cite-dialog}). It shows the recommended citation and a BibTeX
entry, with \gui{Copy Citation} and \gui{Copy BibTeX}\index{citation!BibTeX} buttons that place the text on the
clipboard. The citation is reproduced in Section~\ref{sec:citing}.

\screenshot[0.62\linewidth]{dialogs/citation_dialog}{The \gui{Cite this Software} dialog.}{fig:cite-dialog}{The Cite this Software dialog with the Citation and BibTeX boxes and the Copy Citation, Copy BibTeX and Close buttons.}

\section{Checking for updates}
\label{sec:updates}
\index{updates}

The application checks GitHub for a newer release shortly after it starts and shows the
result in the status bar (\gui{Updates: checking...}, \gui{Updates: up to date}, or a
notice that a new version is available). \menu{About > Check for Updates...}\index{updates!check} repeats the
check on demand:
\begin{itemize}
  \item if a newer version exists, the dialog shows the current and latest version numbers
    with a summary of the release notes, and offers to download the installer for your
    platform to a folder you choose;
  \item if you are up to date, the dialog confirms the installed version;
  \item if the check fails, for example without a network connection, an error message
    explains why.
\end{itemize}
Install the downloaded package as described in Chapter~\ref{ch:installation}.

\section{Reporting a bug}
\label{sec:bug-report}
\index{bug report}\index{diagnostics bundle}

\menu{About > Report a Bug...} opens the bug-report dialog
(Figure~\ref{fig:bug-report}). It contains:
\begin{description}[style=nextline]
  \item[Bug Title] a short summary, for example \emph{Crash when exporting PNG}.
  \item[What happened / Steps to reproduce / Expected behavior] a free-text description,
    pre-filled with a template.
  \item[Generate Diagnostics...] creates a ZIP bundle with a summary of the current session
    (loaded file, plot type, units, color map, combination state and project path),
    environment information (application version, operating system and Python version),
    the operation log, the processing and RFI settings, and a structured provenance
    summary of the loaded data. The path of the bundle is
    shown below the text.
  \item[Copy Issue Text] copies a formatted issue description, including the environment
    summary, to the clipboard.
  \item[Open GitHub Issue] opens a pre-filled issue draft on the project's GitHub page in
    your web browser, where you can attach the diagnostics bundle before submitting.
\end{description}

\screenshot[0.62\linewidth]{dialogs/bug_report_dialog}{The \gui{Report a Bug} dialog.}{fig:bug-report}{The Report a Bug dialog with a title and description filled in, after Generate Diagnostics has created a bundle (its path shown below the text box).}

\begin{note}
  Nothing is sent automatically. The issue is created only when you submit it on GitHub.
  Review the diagnostics bundle before attaching it if your file names or paths are
  confidential.
\end{note}

\section{About}
\label{sec:about}

\menu{About > About} shows the version number, the developer and the affiliation. The
Solar Image Analysis and Maximum Intensities windows have their own \menu{About} menus with
information about those tools.
